
\documentclass{article}
\usepackage{colortbl}
\usepackage{makecell}
\usepackage{multirow}
\usepackage{supertabular}

\begin{document}

\newcounter{utterance}

\centering \large Interaction Transcript for game `wordle\_withclue', experiment `medium\_frequency\_words\_clue\_no\_critic', episode 10 with reflexion{-}Qwen3.5{-}27B.
\vspace{24pt}

{ \footnotesize  \setcounter{utterance}{1}
\setlength{\tabcolsep}{0pt}
\begin{supertabular}{c@{$\;$}|p{.15\linewidth}@{}p{.15\linewidth}p{.15\linewidth}p{.15\linewidth}p{.15\linewidth}p{.15\linewidth}}
    \# & \multicolumn{2}{c}{Player} && \multicolumn{2}{c}{Game Master} \\
    \hline

    \theutterance \stepcounter{utterance}  
    & & & \multicolumn{4}{p{0.6\linewidth}}{
        \cellcolor[rgb]{0.9,0.9,0.9}{
            \makecell[{{p{\linewidth}}}]{
                \texttt{\tiny{[P1$\langle$GM]}}
                \texttt{You have attempted this same task before and did not succeed. These are the plans you wrote after those attempts.} \\
\texttt{Trial \#1: New plan: The previous attempt failed because I violated the instruction "Do not generate any other text" by including a long, rambling internal monologue and logical deduction within the "explanation" field, rather than providing a concise summary. Additionally, I failed to provide a valid "guess:" line with a 5{-}letter word, leaving the response incomplete. Next time, I will strictly adhere to the format by providing only a brief explanation of how the feedback and clue inform the next guess, followed immediately by a valid 5{-}letter word guess. I will perform all logical deductions internally before generating the response to ensure the output contains no extraneous text or incomplete sentences.} \\
\\ 
\\ 
\texttt{You are a language wizard who likes to guess words by using the given rules.} \\
\\ 
\texttt{Welcome to Wordle! You have six attempts to guess the target word, a valid English word of five lowercase letters (a{-}z). Please use the tags "explanation:" and "guess:" to provide a concise explanation for each guess.} \\
\\ 
\texttt{To help you make an informed guess, you will receive a clue for the word, such as} \\
\texttt{clue: snowy white.} \\
\\ 
\texttt{Here is an example guess based on the clue:} \\
\texttt{explanation: In the fairy tail Snow White, the girl is killed because she eats a poisoned apple. And the word apple has 5 letters.} \\
\texttt{guess: apple} \\
\\ 
\texttt{After each guess, your answer will be validated, and you will receive feedback indicating which letters are correct (green), which letters are correct but in the wrong position (yellow), and which letters are incorrect (red). This feedback can be useful in determining which letters to include or exclude in your next guess.} \\
\\ 
\texttt{For example, the feedback for "apple" might be:} \\
\texttt{guess\_feedback: a<yellow> p<yellow> p<green> l<yellow> e<red>} \\
\\ 
\texttt{The explanation should contain details about how both the guess\_feedback and clue work together in tandem to formulate a new guess.} \\
\\ 
\texttt{Do not generate any other text.} \\
\\ 
\texttt{Let's begin with your first guess.} \\
\\ 
\texttt{clue: tithing} \\
            }
        }
    }
     \\ \\

    \theutterance \stepcounter{utterance}  
    & \multicolumn{4}{p{0.6\linewidth}}{
        \cellcolor[rgb]{0.561,0.737,0.561}{
            \makecell[{{p{\linewidth}}}]{
                \texttt{\tiny{[P1$\rangle$GM]}}
                \texttt{explanation: The clue "tithing" refers to the practice of giving a tenth of one's income, often to a church. A common 5{-}letter word associated with this concept is "tithe," which means the act or amount of tithing.} \\
\texttt{guess: tithe} \\
            }
        }
    }
    & & \\ \\

    \theutterance \stepcounter{utterance}  
    & & & \multicolumn{2}{p{0.3\linewidth}}{
        \cellcolor[rgb]{0.9,0.9,0.9}{
            \makecell[{{p{\linewidth}}}]{
                \texttt{\tiny{[GM$|$GM]}}
                \texttt{attempts = 1} \\
\texttt{target\_word = tenth} \\
\texttt{target\_word\_clue = tithing} \\
\texttt{guess = tithe} \\
\texttt{guess\_feedback = t<green> i<red> t<yellow> h<yellow> e<yellow>} \\
            }
        }
    }
    & & \\ \\

    \theutterance \stepcounter{utterance}  
    & & & \multicolumn{4}{p{0.6\linewidth}}{
        \cellcolor[rgb]{0.9,0.9,0.9}{
            \makecell[{{p{\linewidth}}}]{
                \texttt{\tiny{[P1$\langle$GM]}}
                \texttt{guess\_feedback: t<green> i<red> t<yellow> h<yellow> e<yellow>} \\
\\ 
\texttt{Provide your response only in this format.} \\
\texttt{explanation: details} \\
\texttt{guess: word} \\
            }
        }
    }
     \\ \\

    \theutterance \stepcounter{utterance}  
    & \multicolumn{4}{p{0.6\linewidth}}{
        \cellcolor[rgb]{0.561,0.737,0.561}{
            \makecell[{{p{\linewidth}}}]{
                \texttt{\tiny{[P1$\rangle$GM]}}
                \texttt{explanation: The feedback shows 't' is correct in the first position, while 'i', 't', 'h', and 'e' are present but misplaced. Since 't' is already used at the start, the second 't' must be in a different spot. The word "tithe" is a synonym for the clue, but the feedback suggests a different arrangement or a related word. However, "tithe" is the most direct 5{-}letter word for "tithing". Let's re{-}evaluate: if 't' is green at pos 1, and the other letters are yellow, maybe the word is "thief"? No, 'h' and 'e' are yellow in the guess, meaning they are in the word but not at pos 3 and 5. Wait, the clue is "tithing". The word "tithe" fits the clue perfectly. The feedback says 't' (pos 1) is green, 'i' (pos 2) is red (not in word?), 't' (pos 3) is yellow, 'h' (pos 4) is yellow, 'e' (pos 5) is yellow. If 'i' is red, then "tithe" cannot be the word because it contains 'i'. This implies the target word does not contain 'i'. Let's look for a 5{-}letter word related to "tithing" without 'i'. "Tenth" fits "tithing" (a tenth part). Let's check "tenth": t (green), e (yellow in previous guess, so e is in word), n (new), t (yellow in previous guess, so t is in word), h (yellow in previous guess, so h is in word). "Tenth" uses t, e, n, t, h. The previous guess "tithe" had t, i, t, h, e. Feedback: t(green), i(red), t(yellow), h(yellow), e(yellow). So the word has t at pos 1, has t, h, e but not i. "Tenth" has t, e, n, t, h. This matches the letters (t, e, h) and excludes i. The positions: t is 1st. e is 2nd in "tenth", but was 5th in "tithe" (yellow), so e is not 5th. h is 5th in "tenth", but was 4th in "tithe" (yellow), so h is not 4th. t is 4th in "tenth", but was 3rd in "tithe" (yellow), so t is not 3rd. This seems consistent.} \\
\texttt{guess: tenth} \\
            }
        }
    }
    & & \\ \\

    \theutterance \stepcounter{utterance}  
    & & & \multicolumn{2}{p{0.3\linewidth}}{
        \cellcolor[rgb]{0.9,0.9,0.9}{
            \makecell[{{p{\linewidth}}}]{
                \texttt{\tiny{[GM$|$GM]}}
                \texttt{attempts = 2} \\
\texttt{target\_word = tenth} \\
\texttt{target\_word\_clue = tithing} \\
\texttt{guess = tenth} \\
\texttt{guess\_feedback = t<green> e<green> n<green> t<green> h<green>} \\
            }
        }
    }
    & & \\ \\

    \theutterance \stepcounter{utterance}  
    & & & \multicolumn{2}{p{0.3\linewidth}}{
        \cellcolor[rgb]{0.9,0.9,0.9}{
            \makecell[{{p{\linewidth}}}]{
                \texttt{\tiny{[GM$|$GM]}}
                \texttt{game\_result = WIN} \\
            }
        }
    }
    & & \\ \\

\end{supertabular}
}

\end{document}
